\begin{frame}{세 가지 ``왜''}
  \small
  \begin{itemize}
    \item \kb{왜 (1)에서 \texttt{logp\_old}를 저장하는가?}
      $r_t=\pi_\theta/\pi_{\theta_{\mathrm{old}}}$ 의 분모는
      \textbf{데이터를 모았을 때의} 정책 — 그 당시 로그가능도를
      저장해두고, (3)에서 현재 정책의 로그가능도를 계산해 비를
      만든다. 저장하지 않으면 $\theta$ 가 이미 바뀐 뒤 재계산해
      비가 1이 되는 시점이 ``데이터 수집 시점''이 아니라
      ``갱신 시점''이 된다.
    \item \kb{왜 (3)를 여러 에폭 돌리는가?} Week 9.4 재현 버퍼는
      데이터를 \textit{영원히} 저장, PPO는 \textbf{몇 에폭만}
      쓰고 버림 — 클리핑이 ``몇 에폭 재사용''이라는 제한 안에서만
      안정성을 보장.
    \item \kb{왜 (4)에서 버리는가?} on-policy 가정(데이터를 모은
      정책 $\approx$ 현재 정책)이 에폭을 거듭할수록 깨져, $r_t$ 가
      클립 범위를 벗어나는 샘플이 많아지기 때문.
  \end{itemize}
\end{frame}
